% ---- pg-287 （书页 275） ----

MANOVA 结果一致吗? 如果不一致, 你能解释其中的区别吗?

\noindent 6.33\quad 参考习题 6.32, 即

\noindent $X_1=$ 在波长为 560 纳米处光谱反射百分比(绿色)

\noindent $X_2=$ 在波长为 720 纳米处光谱反射百分比(红外线附近)

\hyperref[tab:6.18]{表 6.18} 中的数据是基于上述两个变量的观测, 这些数据是针对三个不同物种(云杉(SS)、日本落叶松(JL)以及黑松(LP))的一岁树苗在成长季节的三个不同时间(150 个阳历日[1]、235 个阳历日[2]和 320 个阳历日[3]). 所有的树苗都在最优营养水平下成长.

\noindent (a) 对\hyperref[tab:6.18]{表 6.18} 中的数据, 完成双因子 MANOVA. 检验品种效应, 时间效应以及物种—时间的交互效应. 取 $\alpha=0.05$.

\noindent (b) 你认为这些数据满足通常的 MANOVA 假设吗? 利用残差分析和对时间的相关观测的可能性进行讨论.

\noindent (c) 林务官对品种和时间的交互作用很感兴趣. 交互作用是不是只对其中一个变量显示出来? 通过进行单变量双因子 ANOVA 分析来对两个响应中的每一个来展开检查.

\noindent (d) 你能想出另外的方法(或者是不同的试验设计)来分析这些数据吗? 这个方法(或者试验设计)要考虑到光谱反射数值的潜在的时间趋势.

\begin{mvtable}
\phantomsection\label{tab:6.18}\mvtabcap{表 6.18\quad 光谱反射数据}
\begin{tabular}{rrccr}
\toprule
560 纳米 & 720 纳米 & 品种 & 时间 & 复制\\\midrule
9.33 & 19.14 & SS & 1 & 1\\
8.74 & 19.55 & SS & 1 & 2\\
9.31 & 19.24 & SS & 1 & 3\\
8.27 & 16.37 & SS & 1 & 4\\
10.22 & 25.00 & SS & 2 & 1\\
10.13 & 25.32 & SS & 2 & 2\\
10.42 & 27.12 & SS & 2 & 3\\
10.62 & 26.28 & SS & 2 & 4\\
15.25 & 38.89 & SS & 3 & 1\\
16.22 & 36.67 & SS & 3 & 2\\
17.24 & 40.74 & SS & 3 & 3\\
12.77 & 67.50 & SS & 3 & 4\\
12.07 & 33.03 & JL & 1 & 1\\
11.03 & 32.37 & JL & 1 & 2\\
12.48 & 31.31 & JL & 1 & 3\\
12.12 & 33.33 & JL & 1 & 4\\
15.38 & 40.00 & JL & 2 & 1\\
14.21 & 40.48 & JL & 2 & 2\\
9.69 & 33.90 & JL & 2 & 3\\
14.35 & 40.15 & JL & 2 & 4\\
38.71 & 77.14 & JL & 3 & 1\\
44.74 & 78.57 & JL & 3 & 2\\
36.67 & 71.43 & JL & 3 & 3\\
37.21 & 45.00 & JL & 3 & 4\\
\bottomrule
\end{tabular}
\end{mvtable}
